%-----------------------------------------------------------------------
\section{Reproducibility and Conclusion}
\label{sec:conclusion}
%-----------------------------------------------------------------------

\subsection{Open Tools and Reproducibility}
Scientific claims require transparent validation. Every line of Python profiling code, raw JSON measurement traces, and algorithmic plotting routines behind the 45,000 data points in this manuscript is strictly open-sourced. We invite aggressive cross-examination of our mathematical constraints. The full benchmark test-suite and compiled data aggregates are permanently mirrored at: \url{https://github.com/[redacted-for-review]}. A permanent DOI will be minted via Zenodo upon acceptance.

\subsection{Conclusion}
We benchmarked five PQC algorithm families — 14 parameter sets total — against the three RFC 7228 device classes, running 45,000 measurement iterations under OS-enforced memory ceilings. We demonstrate statistically significant ($p<0.001$) architectural advantages that favor generic Lattice systems over recursive hash trees across almost all constrained configurations.

Class 0 devices cannot run any current PQC primitive — they require symmetric gateway delegation until lighter-weight post-quantum constructions emerge. For Class 1 networks, ML-KEM-512 paired with FN-DSA-512 (hardware FPU mandatory) offers the best balance of speed, memory, and transmission overhead. Class 2 hardware comfortably supports ML-KEM-768 and ML-DSA-2/3, including hybrid classical modes on Ethernet-class links.

The energy model ($E_{total} = E_{cpu} + E_{tx}$) confirms that ML-KEM-512's fast computation offsets its larger payload: total energy on Class 2 is 3.8 mJ, competitive with classical ECDH. The TOPSIS rankings, ANOVA validation, and feasibility matrices presented here give IoT architects a concrete, data-backed selection guide — not a general recommendation, but specific algorithm-to-hardware mappings with quantified trade-offs.
